\documentclass[11pt]{article}

\usepackage[T1]{fontenc}
\usepackage{amsmath,amssymb,mathtools}
\usepackage{booktabs,longtable,tabularx,array}
\usepackage[margin=1in]{geometry}
\usepackage{microtype}
\usepackage{enumitem}
\usepackage{xcolor}
\usepackage{hyperref}
\usepackage{titlesec}

\definecolor{irisblue}{HTML}{22577A}
\definecolor{irisgray}{HTML}{4B5563}
\hypersetup{
  colorlinks=true,
  linkcolor=irisblue,
  urlcolor=irisblue,
  pdftitle={An Endpoint-Local Adaptive L2 Rule for the Full-HLA Model},
  pdfauthor={IRIS methodological development memo}
}
\titleformat{\section}{\Large\bfseries\color{irisblue}}{\thesection}{0.65em}{}
\titleformat{\subsection}{\large\bfseries\color{irisgray}}{\thesubsection}{0.65em}{}
\setlength{\parindent}{0pt}
\setlength{\parskip}{0.65em}
\setlist{itemsep=0.25em,topsep=0.35em}

\newcommand{\sigmoid}{\operatorname{sigmoid}}
\newcommand{\LSE}{\operatorname{LSE}}
\newcommand{\AP}{\operatorname{AP}}

\title{\textbf{An Endpoint-Local Adaptive Level~2 Rule for the Full-HLA Model:\\Distinct-Epitope Breadth with HLA Corroboration}}
\author{IRIS methodological development memo}
\date{September 5, 2026}

\begin{document}
\maketitle

\begin{abstract}
The complete-$F$ rerun showed that the original adaptive Level~2 family can
improve on fixed maximum within each of PDAC, SARS-CoV-2 SPIKE, and
SARS-CoV-2 NONSPIKE, but no single original parameter tuple transported across
all three. A cohort-relative gate passed the numerical evidence criterion, but
was scientifically inadmissible because the score of one long peptide depended
on the other peptides in the scoring batch. A sequence of endpoint-local
alternatives then isolated two complementary kinds of evidence. First,
candidate breadth should be counted across distinct short peptides rather than
across peptide--HLA tuples. Second, support through a second distinct HLA is
useful, but not strong enough to replace epitope-level aggregation by itself.
Combining these ideas yields a deterministic score computed only from one long
peptide's candidate $F$ values and HLA identities. A representative central
tuple passes the staged paired-CNAP criterion against fixed maximum in all
three datasets at 600 replications, and 61 of 81 nearby tuples do likewise.
These are encouraging post-result model-development findings, not independent
validation and not yet a production replacement.
\end{abstract}

\section{The question that Level~2 must answer}

The assay label belongs to a long peptide, whereas the Full-HLA Level~1 model
assigns a score to each candidate short-peptide--HLA pair. Level~2 must turn
that roster into one score that ranks the long peptide. Let
\begin{equation}
z_{jh}=\log(F_{jh}+\epsilon), \qquad \epsilon=10^{-12},
\end{equation}
where $j$ indexes a distinct candidate 9--12-mer and $h$ indexes an HLA allele
in the endpoint's profile. The desired aggregate must satisfy a strict locality
condition:
\begin{quote}
The score of a long peptide may depend on that peptide, its candidate Level~1
scores, its own HLA profile, and globally frozen constants. It may not depend
on the identities or scores of other long peptides in the dataset.
\end{quote}
This is both a deployment requirement and a biological one. Adding an
unrelated peptide to an evaluation cohort should not change the meaning of the
peptide already being scored.

\section{The original adaptive rule, briefly}

The adaptive Hill rule is developed in detail elsewhere; only the pieces
needed for the present argument are repeated here. For a roster of log scores
$z_1,\ldots,z_K$, its power-mean anchor is
\begin{equation}
A_\alpha(z)=
\begin{cases}
K^{-1}\sum_i z_i, & \alpha=0,\\[3pt]
\alpha^{-1}\log\!\left(K^{-1}\sum_i e^{\alpha z_i}\right),
  & 0<\alpha<\infty,\\[3pt]
\max_i z_i, & \alpha=\infty.
\end{cases}
\end{equation}
Thus $\alpha$ moves the anchor from the mean of log scores toward maximum.
After removal of the numerical floor, define
\begin{equation}
x_i=\max(e^{z_i}-\epsilon,0),\qquad
p_i=\frac{x_i}{\sum_r x_r},\qquad
U=\log\!\left(\epsilon+\sum_i x_i\right).
\end{equation}
The Hill order $q$ measures how broadly the positive signal is distributed.
For the ultimately useful case $q=\infty$,
\begin{equation}
D_\infty=1-\max_i p_i.
\end{equation}
The available corroboration offer is
\begin{equation}
C_{\alpha,q}=D_q\max(U-A_\alpha,0),
\end{equation}
and the adaptive score is the unique solution of
\begin{equation}
S=A_\alpha+C_{\alpha,q}
\sigmoid\!\left(\frac{c-S}{\kappa}\right).
\label{eq:original-adaptive}
\end{equation}
The sigmoid is a deterministic smooth switch, not a response probability.
The center $c$ says where corroboration begins to close, while $\kappa$
controls the sharpness of that transition.

\section{What the complete-$F$ results changed}

The complete-$F$ selection run removed the ambiguity caused by uncomputed
candidate tensors. It then exposed a transport problem in the original
Full-HLA PR rule. Finite-$\alpha$ anchors were vulnerable to dilution by the
large complete candidate roster. At $\alpha=\infty$, however, useful
corroboration was present in every dataset. The best supported absolute gate
centers were nevertheless different: approximately $-1.2$ for PDAC, $-3.8$
for SPIKE, and $-2.0$ for NONSPIKE. No common tuple from the original frozen
family passed the all-context eligibility requirement.

This was not evidence that maximum was unbeatable. Rather, maximum supplied a
transportable anchor while the useful location of the corroboration gate did
not transport. The subsequent investigation asked what endpoint-intrinsic
structure could identify useful corroboration without letting the cohort set
the score.

\subsection{A successful diagnostic that could not be the answer}

Replacing the absolute gate center by the 38th percentile of the evaluation
batch's maximum scores produced positive supported magnitudes in all three
datasets at 600 replications. This established that the original breadth offer
contained predictive information when admitted in the appropriate score
region. It did not establish an admissible scoring rule. If
\begin{equation}
c_E=Q_{0.38}\{m_i:i\in E\},
\end{equation}
then changing the other members of evaluation set $E$ can change the score of
endpoint $i$. The rule was therefore discarded as transductive. Its value was
diagnostic: it showed that there was secondary signal to recover, and that the
problem was how that signal was represented and gated.

\subsection{The key distinction: two kinds of breadth}

The raw Full-HLA roster confounds two different claims:
\begin{enumerate}
  \item several \emph{distinct short peptides} within the long peptide have
  meaningful Level~1 signal; and
  \item meaningful signal is available through more than one \emph{distinct
  HLA} in the endpoint's HLA profile.
\end{enumerate}
A short peptide that scores under several HLAs should not automatically count
as several separate epitopes. Conversely, collapsing everything to one score
per short peptide can hide the fact that the patient has more than one credible
presentation route. The successful construction treats these as separate,
complementary features.

\section{The endpoint-local hybrid rule}

\subsection{Step 1: collapse to distinct short peptides}

For each distinct short peptide $j$, retain its strongest HLA-specific score:
\begin{equation}
u_j=\max_h z_{jh}.
\label{eq:peptide-collapse}
\end{equation}
The epitope roster is $\mathcal U=(u_1,\ldots,u_J)$. Its maximum
\begin{equation}
m=\max_j u_j=\max_{j,h}z_{jh}
\end{equation}
is exactly the ordinary Full-HLA maximum. The collapse changes the accounting
of breadth, not the identity of the strongest candidate.

\subsection{Step 2: compute adaptive distinct-epitope breadth}

On $\mathcal U$, define
\begin{align}
x_j&=\max(e^{u_j}-\epsilon,0),\\
p_j&=x_j\Big/\sum_r x_r,\\
D&=1-\max_j p_j,\\
U&=\log\!\left(\epsilon+\sum_jx_j\right),\\
C&=D\max(U-m,0).
\end{align}
The distinct-epitope score $P$ is the unique solution of
\begin{equation}
P=m+C\sigmoid\!\left(\frac{-2.2-P}{0.13}\right).
\label{eq:peptide-branch}
\end{equation}
This is Equation~\eqref{eq:original-adaptive} with
$\alpha=\infty$, $q=\infty$, $c=-2.2$, and $\kappa=0.13$, applied after
short-peptide collapse.

\subsection{Step 3: measure support from a second distinct HLA}

Collapse in the other direction to obtain one maximum for each HLA:
\begin{equation}
M_h=\max_j z_{jh}.
\end{equation}
Let $m_2$ be the second-largest value among the distinct-HLA maxima $M_h$.
Define a bounded HLA corroboration branch
\begin{equation}
H=m+B\sigmoid\!\left(\frac{m_2-t}{\delta}\right).
\label{eq:hla-branch}
\end{equation}
For the representative central tuple,
\begin{equation}
t=-6.45,\qquad \delta=0.02,\qquad B=1.
\end{equation}
This branch asks a narrow question: is there at least one credible candidate
route through a second HLA? It does not average the two best HLAs and does not
allow the second HLA to replace the best candidate.

\subsection{Step 4: combine the two branches}

The final score is the convex combination
\begin{equation}
\boxed{
S_{\mathrm{hybrid}}=(1-w)P+wH,
\qquad w=0.12.
}
\label{eq:hybrid}
\end{equation}
Equivalently, 88\% of the score comes from maximum-anchored adaptive breadth
among distinct short peptides and 12\% from the second-HLA branch. Because the
HLA bonus is itself bounded by $B$, its maximum direct contribution to the
final score is $wB=0.12$ log-score units. The second HLA can refine rankings;
it cannot dominate them.

\subsection{A literal computation recipe}

For one endpoint, the calculation is:
\begin{enumerate}
  \item read its aligned $(\text{short peptide},\text{HLA},z)$ roster;
  \item group by short-peptide sequence and retain the maximum $z$ within each
  group;
  \item compute $m$, $D$, $U$, and $C$ from this collapsed roster;
  \item solve Equation~\eqref{eq:peptide-branch} for $P$ by bisection;
  \item group the original roster by HLA and retain the maximum $z$ within
  each group;
  \item take the second-largest HLA maximum $m_2$ and compute $H$ from
  Equation~\eqref{eq:hla-branch};
  \item combine $P$ and $H$ using Equation~\eqref{eq:hybrid}.
\end{enumerate}
Every input in this recipe belongs to the endpoint being scored. All other
quantities are frozen constants.

\section{Results}

The representative central tuple was compared with the same Full-HLA fixed
maximum reference using the staged paired-CNAP procedure with 600
replications. ``Supported magnitude'' is the conservative retained improvement
used by the eligibility rule; ``survival'' is the fraction of the staged
replication subsets retaining the required direction.

\begin{center}
\small
\begin{tabular}{@{}lrrr@{}}
\toprule
Dataset & Observed retained difference & Supported magnitude & Survival \\
\midrule
PDAC      & $+0.0003428$ & $+0.0001200$ & $0.9344$ \\
SPIKE     & $+0.0000340$ & $+0.00000589$ & $0.9280$ \\
NONSPIKE  & $+0.0007612$ & $+0.0001587$ & $0.9933$ \\
\bottomrule
\end{tabular}
\end{center}

All three supported magnitudes are positive and all three survival fractions
exceed the declared $0.81$ threshold. SPIKE remains the limiting dataset by
supported magnitude, but it is no longer a failure.

The result is not confined to one exact numerical point. A 600-replication
neighborhood varied
\begin{align*}
t&\in\{-6.475,-6.45,-6.425\},\\
\delta&\in\{0.015,0.020,0.025\},\\
w&\in\{0.08,0.10,0.12\},\\
B&\in\{0.50,0.75,1.00\}.
\end{align*}
Of the $3^4=81$ combinations, 61 passed in all three datasets. Passing tuples
occurred at every tested value of each parameter. This neighborhood result is
stronger evidence than a single grid hit, although it does not remove the
post-selection status of the entire exercise.

\section{High-level interpretation}

The hybrid is most naturally read as a decomposition of corroboration rather
than as a statistical probability model.

\textbf{The maximum remains the sufficiency anchor.}
The strongest predicted short-peptide--HLA route is protected. The successful
rule did not require a finite-$\alpha$ compromise below maximum.

\textbf{Epitope breadth is counted in biological units.}
The Hill calculation is performed across distinct short-peptide sequences.
This prevents the same peptide from manufacturing apparent breadth merely
because it is represented for several HLAs.

\textbf{HLA breadth is separate corroboration.}
After peptide collapse, the second-best distinct HLA still contains useful
information. A credible second HLA suggests that the long peptide has more
than one presentation route in that patient's actual HLA profile. The small
mixture weight reflects the evidence: this feature is useful as a modifier,
not as the primary aggregate.

\textbf{The two gates answer different questions.}
The adaptive gate in Equation~\eqref{eq:peptide-branch} regulates how much
multi-epitope accumulation a long peptide receives. The second-HLA switch in
Equation~\eqref{eq:hla-branch} recognizes an additional presentation route.
Neither sigmoid is a probability. Both are smooth deterministic versions of
threshold decisions.

\textbf{The rule is endpoint-local.}
Long-peptide sequence determines the candidate short peptides; the endpoint's
HLA profile determines their presentation routes; the complete Level~1 model
determines the $F$ scores. No cohort name, label prevalence, batch quantile,
other endpoint, decoy distribution, or fitted endpoint-specific random
quantity enters the calculation.

\section{What was tried and did not provide the final rule}

The negative results are part of the argument. They show that the successful
hybrid is not simply ``use the second HLA'' or ``estimate a personalized gate
center.'' The catalog below records distinct functional forms, rather than
every parameter value evaluated within each form.

\subsection{Original and location-adaptive gates}

\begin{longtable}{@{}p{0.27\textwidth}p{0.34\textwidth}p{0.31\textwidth}@{}}
\toprule
Functional form & Intended interpretation & Result and lesson \\
\midrule
\endfirsthead
\toprule
Functional form & Intended interpretation & Result and lesson \\
\midrule
\endhead
Raw-roster self-gated Hill rule,
$S=A_\alpha+C_{\alpha,q}\sigmoid((c-S)/\kappa)$
& Allow the candidate roster to choose a mean-to-maximum anchor and a
corroboration increment.
& No common Full-HLA PR tuple passed all three datasets. Finite $\alpha$ was
particularly vulnerable to dilution; useful $\alpha=\infty$ regions required
different absolute $c$ values. \\
\addlinespace
Batch-quantile center,
$c_E=Q_\rho\{m_i:i\in E\}+\Delta$
& Put the gate at the same relative location in each evaluation population.
& Numerically successful at $\rho=0.38$, but rejected as transductive: the
score of one long peptide changes when other peptides enter or leave $E$.
This was a diagnostic success and a scientific failure. \\
\addlinespace
Candidate-count center,
$c_i=a+b\log K_i$
& Adjust the gate for the number of enumerated opportunities.
& Reduced ordinary AP substantially in PDAC and slightly in SPIKE. Raw count
does not encode redundancy, HLA composition, or upper-tail geometry. \\
\addlinespace
Raw-roster location centers,
$c_i=T(z_i)+\Delta$, with $T$ equal to mean, median, 75th, 90th, or 95th
percentile
& Make $c$ depend on the endpoint's own candidate-score distribution.
& Some observed improvements occurred, especially for the upper quartile, but
no common staged-supported rule emerged. One location statistic did not
separate useful from false corroboration. \\
\addlinespace
HLA-location centers: median or second-largest of $\{M_h\}$, plus an offset
& Calibrate the gate using the endpoint's own HLA-specific maxima.
& Did not yield an all-context supported tuple. HLA location alone was
insufficient. \\
\addlinespace
Generalized intrinsic centers: a power mean of the raw roster, HLA maxima,
short-peptide maxima, or within-peptide second-HLA scores, plus an offset
& Let several endpoint-intrinsic roster summaries define $c_i$.
& Observed screens produced candidates, but the compact paired-CNAP tests did
not support a common rule. Greater flexibility in the center did not solve the
representation problem. \\
\addlinespace
HLA-dominance gate,
$S=m+C\sigmoid((g-\tau)/\delta)$, where $g=M_{(1)}-M_{(2)}$
& Admit breadth according to whether one HLA dominates the rest; both the
full-roster and winning-HLA offers were tested.
& None of 32 tuples passed all three datasets at 600 replications. A gap can
confound one strong route with two weak routes and discards the absolute
strength of the second HLA. \\
\addlinespace
Ungated constant fraction,
$S=m+\lambda C$
& Avoid choosing a gate location entirely.
& No common positive ordinary-AP improvement. Breadth was not globally
positive-label enriched and needed selective admission. \\
\bottomrule
\end{longtable}

\subsection{Direct top-two-HLA functions}

Let $M_{(1)}\geq M_{(2)}$ be the two largest maxima from distinct HLAs. These
experiments established that $M_{(2)}$ contains signal, but also that it is not
a sufficient Level~2 summary by itself.

\begin{longtable}{@{}p{0.34\textwidth}p{0.26\textwidth}p{0.32\textwidth}@{}}
\toprule
Functional form & Intended interpretation & Result and lesson \\
\midrule
\endfirsthead
\toprule
Functional form & Intended interpretation & Result and lesson \\
\midrule
\endhead
Power mean of $M_{(1)}$ and $M_{(2)}$
& Treat the two best HLA routes as the entire roster.
& No common supported rule. Averaging can penalize a valid leader and ignores
distinct-epitope breadth. \\
\addlinespace
Weighted log-score mean,
$(1-r)M_{(1)}+rM_{(2)}$
& Give the second HLA a fixed fractional vote.
& No common supported rule; the second HLA should modify rather than dilute
the maximum. \\
\addlinespace
Exponential agreement bonus,
$M_{(1)}+B\exp[-(M_{(1)}-M_{(2)})/s]$
& Reward two HLA maxima that are close together.
& No common supported rule. Closeness does not distinguish two strong routes
from two weak ones. \\
\addlinespace
Hinge agreement bonus,
$M_{(1)}+r[d-(M_{(1)}-M_{(2)})]_+$
& Give a linear bonus when the two HLA routes are sufficiently close.
& No common supported rule for the same reason: relative agreement lacks an
absolute evidence requirement. \\
\addlinespace
Thresholded second-HLA bonus,
$M_{(1)}+B\sigmoid((M_{(2)}-t)/\delta)$
& Reward a second HLA only when its own score is strong enough.
& This was the most promising HLA-only form, but no common standalone tuple
was stable across PDAC, SPIKE, and NONSPIKE. It survives in the final rule only
as a low-weight branch. \\
\addlinespace
Length-adjusted threshold,
$t_i=t_0+s(L_i-15)_+$
& Allow longer peptides, with more candidate windows, to require stronger
second-HLA evidence.
& Did not restore common support. Length was too coarse a description of the
actual peptide--HLA opportunity structure. \\
\bottomrule
\end{longtable}

\subsection{Adaptive top-two and structural functions}

The top-two idea was also inserted into the adaptive rule in several ways.
Offers were computed from only the two HLA maxima, from all candidates
belonging to the two leading HLAs, and from all raw candidates. Anchors included
finite power means and maximum. Gate centers included a common absolute $c$,
$M_{(2)}+\Delta$, $aM_{(2)}+b$, and
$(M_{(1)}+M_{(2)})/2+\Delta$. None produced one staged-supported tuple across
all three datasets. These failures showed that moving the existing gate center
or restricting its offer roster did not cleanly separate epitope multiplicity
from HLA multiplicity.

More literal structural rules were then tested:
\begin{longtable}{@{}p{0.37\textwidth}p{0.25\textwidth}p{0.30\textwidth}@{}}
\toprule
Functional form & Intended interpretation & Result and lesson \\
\midrule
\endfirsthead
\toprule
Functional form & Intended interpretation & Result and lesson \\
\midrule
\endhead
Power mean of the best two HLA scores for the \emph{same} short peptide
& Require promiscuous presentation of one exact epitope.
& Did not transport. The requirement was too restrictive and discarded
evidence from different epitopes presented by different HLAs. \\
\addlinespace
$\max\{m,M_{(2)}+b\}$
& Let a sufficiently strong second HLA lift the maximum.
& No common supported rule as a standalone score. \\
\addlinespace
$\max\{m,u_{(2)}+b\}$, where $u_{(2)}$ is the second distinct-peptide maximum
& Require a second credible epitope.
& Improved selected settings but did not satisfy the joint criterion. A hard
two-epitope summary discarded the remainder of the breadth geometry. \\
\addlinespace
$\max\{m,v_{(2)}+b\}$, where $v_{(2)}$ is the second-HLA score for one exact
short peptide
& Reward one peptide supported by two HLAs.
& Did not transport; again too restrictive. \\
\addlinespace
$\max\{m,(v_{(1)}+v_{(2)})/2+b\}$
& Use the mean of two HLA scores for one peptide as corroboration.
& Did not transport and reintroduced dilution within the corroborating branch. \\
\addlinespace
$m+r\min\{[d-(m-M_{(2)})]_+,[d-(m-u_{(2)})]_+\}$
& Require HLA breadth and distinct-peptide breadth simultaneously.
& No common supported rule. A hard conjunction was less effective than
computing epitope breadth continuously and using HLA evidence as a modifier. \\
\bottomrule
\end{longtable}

\subsection{Collapsed-roster and hybrid functions}

Applying the adaptive Hill rule after collapsing to one maximum per HLA, to
the top two HLA maxima, or to one maximum per distinct short peptide directly
tested the unit in which breadth should be counted. HLA-maxima and top-two-HLA
rosters did not yield a common rule. Distinct-short-peptide collapse was the
decisive near-success. With $\alpha=q=\infty$, $c=-2.2$, and
$\kappa=0.13$, it passed PDAC and NONSPIKE at 600 replications but missed
SPIKE only through a slightly negative supported magnitude:
\begin{center}
\small
\begin{tabular}{@{}lrr@{}}
\toprule
Dataset & Supported magnitude & Survival \\
\midrule
PDAC     & $+0.0000254$ & $0.8493$ \\
SPIKE    & $-0.00000564$ & $0.9024$ \\
NONSPIKE & $+0.0001617$ & $0.9933$ \\
\bottomrule
\end{tabular}
\end{center}
This result located the missing information: the epitope representation was
substantially right, but SPIKE still benefited from a small HLA-specific cue.

The first hybrid moved the peptide-collapsed score a gated fraction toward its
collapsed log-sum-exp ceiling,
\begin{equation}
P+wG_{\mathrm{HLA}}(U-P).
\end{equation}
That bounded form did not provide the common supported rule. The successful
form instead constructs the simple HLA branch $H=m+BG_{\mathrm{HLA}}$ and
takes a convex combination with $P$. This keeps the two biological claims
legible: continuous breadth across distinct epitopes in one branch and bounded
evidence from a second presentation route in the other.

\section{Evidential status and next decision}

The numerical claim is verified on the present complete-$F$ PDAC, SPIKE, and
NONSPIKE inputs under the same staged paired-CNAP machinery used by the
selector. The conceptual claim is narrower: separating distinct-epitope breadth
from second-HLA corroboration produces a locally robust, endpoint-local Full-HLA
score that outperforms fixed maximum by the declared selection criterion in all
three datasets.

Three limitations should remain attached to that statement:
\begin{enumerate}
  \item The rule and its constants were developed after inspecting these
  results. Re-evaluation on the same datasets is post-selection evidence, even
  with 600 replications.
  \item The PR improvements are small in absolute units, especially for SPIKE.
  Their importance lies in consistent direction and resampling support, not in
  a large change in empirical AP.
  \item The exploratory implementation has not replaced the production
  \texttt{select\_adaptive\_hillq} contract. The rule should be reviewed,
  frozen, implemented, and tested before any new tournament; genuinely held-out
  data would provide the cleanest confirmation.
\end{enumerate}

The appropriate conclusion is therefore neither that the search failed nor
that a definitive production model has been proven. The search isolated a
specific biological decomposition, found a simple computable rule expressing
it, and verified a broad local region of numerical support. The next scientific
act is to freeze that proposal before asking it another question.

\section*{Reproducibility record}

The representative and neighborhood results are stored in:
\begin{itemize}
  \item \path{artifacts/iris_fullroster_complete_f_20260904/adaptive_rule_exploration/full_hla_hybrid_convex_cnap_600.json};
  \item \path{artifacts/iris_fullroster_complete_f_20260904/adaptive_rule_exploration/full_hla_hybrid_convex_robustness_cnap_600.json}.
\end{itemize}
The exploratory implementation is
\path{supported_ap_code/crates/select_adaptive_hillq/src/bin/explore_full_hla_hybrid_breadth.rs}.
The original rule is documented in
\path{supported_ap_code/crates/select_adaptive_hillq/popperian_adaptive_l2_recommendation.tex}.

\end{document}
